\section{What the Presence Head Responds To}
\label{sec:mechanism}

\begin{table}[t]
    \centering
    \small
    \setlength{\tabcolsep}{3.2pt}
    \resizebox{\columnwidth}{!}{\begin{tabular}{@{}lrcccc@{}}
\toprule
Prompt & $n$ & $s_p\!\geq\!0$ & mean $s_p$ & any mask & object \\
\midrule
\multicolumn{6}{@{}l}{\emph{Absent phrase, LVIS held-out}} \\
text only & 1,500 & 1.6 & $-5.69$ & 1.3 & -- \\
+ box on largest object & 1,500 & 99.5 & $9.17$ & 99.5 & 96.8 \\
+ box on background & 500 & 99.6 & $8.74$ & 99.6 & -- \\
\multicolumn{6}{@{}l}{\emph{Absent phrase, LVIS selection images}} \\
text only & 400 & 2.0 & $-5.73$ & -- & -- \\
+ 1 box & 400 & 99.8 & $9.29$ & 99.8 & 96.3 \\
+ 2 boxes & 305 & 100.0 & $9.96$ & 100.0 & 94.4 \\
+ 3 boxes & 226 & 100.0 & $10.00$ & 100.0 & 95.1 \\
+ box at half scale & 400 & 100.0 & $9.33$ & 100.0 & 0.3 \\
+ box at double scale & 400 & 99.8 & $8.83$ & 99.8 & 17.8 \\
+ box shifted by half its width & 400 & 100.0 & $9.15$ & 100.0 & 3.8 \\
\multicolumn{6}{@{}l}{\emph{COCO (absence by the 80-class label set)}} \\
absent phrase + \emph{negative} box & 1,000 & 0.0 & $-9.88$ & 0.0 & -- \\
positive box, no text & 500 & 100.0 & $9.97$ & 100.0 & 90.8 \\
other present class + box & 421 & 100.0 & -- & -- & 86.5 \\
\bottomrule
\end{tabular}
}
    \caption{\textbf{Controls.}
    Percent of episodes, and the mean presence logit.
    \emph{Object}: a kept mask covers the anchor.
    Boxes are positive exemplars unless marked.
    Rows with two and three boxes use the images that have that many further instances.}
    \label{tab:controls}
\end{table}

A box could raise presence for a reason that has something to do with the image: there is an object in it, and the head might take the phrase as a loose name for that object.
The controls in \cref{tab:controls} rule this out.

\paragraph{Background.}
We place a box of the anchor's size at the position, on a $12 \times 12$ grid, where it overlaps the fewest annotated pixels, and prompt with a verified-absent phrase.
On the held-out images presence is non-negative in 99.6\% of 500 episodes, with a mean logit of 8.74, and a mask is returned in 99.6\%.
On the selection images the returned mask covers the box region in 37.4\% of episodes; elsewhere the model finds something else to return.
LVIS cannot certify that such a region is empty, only that nothing in it was annotated.

\paragraph{Perturbed and multiple boxes.}
On 400 selection images we shrink the anchor box to half its size, enlarge it to double, or shift it by half its width.
Presence stays non-negative in 100\%, 99.8\% and 100\% of episodes, while a kept mask covers the anchor in 0.3\%, 17.8\% and 3.8\%.
The head does not need the box to hold the object, and the masks follow the box.
With two and three positive boxes the rate is 100\% and the mean logit approaches the clamp (9.96 and 10.00).

\paragraph{The label decides.}
Flipping the label reverses the effect.
With the anchor box as a lone \emph{negative} exemplar and an absent phrase, presence is negative in all 1,000 COCO episodes, at the lower clamp of $-10$ in 95.4\%, and no mask is returned.
With a positive box and no text at all, the mean logit is 9.97.

\paragraph{An account from the training protocol.}
The \sam{} paper introduces the presence token in the training stage that also introduces interactive exemplars.
In that stage, positive boxes are the ground-truth boxes of instances the model missed, and negative boxes are confident false positives without significant overlap with ground truth~\cite{carion2026sam3}.
A missed instance exists only when the concept is present.
So, in the protocol as published, a positive exemplar is not paired with an absent concept, and its label alone predicts presence.
A false positive can occur whether or not the concept is present, so a negative label carries no such guarantee.
Before the runs we wrote down two predictions from this account: a positive box on background would still pin presence high, and a lone negative box with an absent phrase would pin it low.
Both hold in \cref{tab:controls}.
The negative box in fact pushes presence lower than the phrase alone does (mean $-9.88$ against $-6.21$), which would follow if prompts whose only exemplars are negative arise mostly on images where every detection is a false positive.
The published evaluation could not have exposed the problem, because exemplar prompts were scored only on images where the prompted object appears~\cite{carion2026sam3}.
The account is consistent with what we measured, and it is the textbook form of a shortcut~\cite{geirhos2020shortcut}.
It remains an account.
Confirming it would take a training run in which positive exemplars are paired with absent phrases, and we have not done one.
